\documentclass[11pt]{article}
\usepackage[margin=1in]{geometry}
\usepackage{amsmath,amssymb,amsthm,booktabs}
\usepackage[hidelinks]{hyperref}
\setlength{\emergencystretch}{2em}
\newtheorem{theorem}{Theorem}
\newtheorem{proposition}[theorem]{Proposition}
\newcommand{\mr}{\operatorname{mr}}
\title{Double-transposition quotients: a sharp known control and an exact boundary}
\author{Internal research working proof}
\date{17 November 2026}
\begin{document}
\maketitle

The retained sixteen targets have no transitive degree-eight or
degree-nine permutation representation sending every meridian to
two disjoint transpositions. These are exactly the transitive
degrees where the elementary permutation-index bound below could
reach their certified bridge upper bound four. This is a finite
route exclusion; the sixteen knots and full MRC remain open.
The test succeeds on known torus controls, including a full
$A_8$ image with a meridian class distinct from binary
four-dimensional transvections.

\section{The elementary bound and its precise scope}

For a permutation $g\in S_r$, let $\delta(g)$ be $r$ minus its
number of cycles, including fixed points. This is its minimum
transposition length.
\begin{proposition}
If $\phi:\pi_1(S^3\setminus K)\to S_r$ is transitive and a
meridian has permutation index $\delta>0$, then
\[
                 \mr(K)\ge\left\lceil\frac{r-1}{\delta}\right\rceil.
\]
\end{proposition}
\begin{proof}
Take any meridional generating family of size $d$. Its images
are conjugates or inverses of the chosen meridian image, so all
have index $\delta$. Factor each image into exactly $\delta$
transpositions, and place all those edges on the $r$ sheets.
Every generating permutation preserves every component of this
graph. Transitivity therefore forces it to be connected. A
connected graph on $r$ vertices has at least $r-1$ edges, while
this graph has at most $d\delta$. Thus $r-1\le d\delta$.
\end{proof}

This is the standard transposition-graph argument extended to
arbitrary cycle type; no priority claim is made. Earlier
meridian-preserving symmetric and alternating quotient
constructions, including finite-cycle meridians, are credited
in \cite{bk,jp}. When meridians are double transpositions,
$\delta=2$. A four-meridian upper certificate forces $r\le9$,
and the displayed lower bound reaches four exactly at $r=8,9$.
We exclude those two degrees only on the unchanged sixteen
PDs. Lower-degree representations could still support an
enhanced cover-rank or other bound, and meridians with three
or more transpositions are not covered by this exclusion.

\section{A full alternating-group control on known knots}

Permutations act on $\{0,\ldots,7\}$, with the right factor
applied first. Set
\[
 X=(0\ 7\ 5\ 1)(2\ 6\ 4\ 3),\quad
 Y=(0\ 1\ 2\ 3\ 4)(5\ 6\ 7),\quad
 M=(0\ 6)(2\ 5).
\]
Direct evaluation gives $X^4=Y^{15}=1$ and $XM=Y^4$.
Exact finite enumeration gives $|\langle X,Y\rangle|=20160$.
Both generators are even, so this group is all of $A_8$.
The independent checker enumerates by left multiplication,
separately from the producer's right multiplication.

\begin{theorem}
For every $j\ge0$, $T(4,60j+15)$ admits a quotient onto $A_8$
whose meridians are double transpositions. The elementary
bound and four-strand braid certify $\mr=b=4$.
\end{theorem}
\begin{proof}
Put $q=60j+15$. In the standard torus presentation and meridian
convention retained from \cite[Section 5.1]{cornwell},
\[
 \Gamma=\langle x,y\mid x^4=y^q\rangle,\qquad
 m=x^{-1}y^{15j+4}.
\]
The exponents satisfy $-q+4(15j+4)=1$. Sending $x,y$ to $X,Y$
satisfies the relation and sends $m$ to $X^{-1}Y^4=M$.
The image is the same full $A_8$ for every $j$. Its degree
eight and index-two meridians give $\mr\ge4$. The closure of
the four-strand torus braid gives $b\le4$.
\end{proof}

These are already settled torus knots; this explicit quotient
does not claim a new MRC case or worldwide novelty. The
classical torus result is due to Rost--Zieschang \cite{rz}.

This full image lies outside binary three-dimensional affine
images, since $|A_8|=20160>1344=|\operatorname{AGL}_3(2)|$.
Although $A_8\cong\operatorname{GL}_4(2)$ \cite{atlas}, its
double-transposition class has size $210$: there are
$3\binom84$ such elements, and an odd permutation on fixed
points centralizes a representative, so its $S_8$ class does
not split in $A_8$. In contrast a binary transvection
$I+vf$ has $v\ne0$, $f\ne0$, $f(v)=0$, giving $15\cdot7=105$
distinct matrices. All are conjugate by a basis change.
An isomorphism cannot exchange classes of unequal size.
Thus the pinned external binary rank-one search does not
include this meridian class. This statement concerns search
coverage, not any of its valid deposited certificates.
No classification of higher-field representations is inferred.

\section{Exact finite coverage and independent replay}

The two plans were frozen separately before execution: first
degree eight, then the remaining degree nine. Each uses the
original sixteen PDs and retained four-seed traces. All ordered
four-tuples are considered up to simultaneous $S_r$ conjugacy.
Conjugate the first entry to $(0\ 1)(2\ 3)$; its centralizer has
order $8(r-4)!$. Successive stabilizers select canonical tail
representatives. Only tuples whose transposition-edge graph
is connected can generate a transitive image.

\begin{center}
\begin{tabular}{rrrrr}
\toprule
$r$ & Class size & Raw tails & All tuple orbits & Transitive orbits\\
\midrule
8 & 210 & 9,261,000 & 58,429 & 18,450\\
9 & 378 & 54,010,152 & 86,231 & 6,561\\
\bottomrule
\end{tabular}
\end{center}

No transitive representative extends to a valid coloring on
any target, in either degree. Each successful coloring would
be determined by the four seed meridians, so this covers every
representation of the stated type, including proper images.
No classification of transitive subgroups is assumed.

The independent verifier enumerates the meridian class by
filtering all permutations, rather than by choosing supports
and pairings. It computes full centralizer orbits on the raw
tail cube and marks every tuple, checking disjointness and
that no tuple is omitted. This does not rely only on an
orbit-count sum or the producer's prefix recursion.

Geometric paths between undercrossings independently rebuild
the arcs and check the legacy index transport. The source PD
parser is reused only for those old index names. The verifier
checks each upper trace, assigns all remaining images in one
pass through it, and then tests every crossing pointwise.
The producer instead uses iterative propagation with early
failure. Both obtain zero hits on all sixteen targets.

The known $T(4,7)$ control yields fourteen transitive pairs of
image order $1344$; $T(4,15)$ yields fifteen of order $20160$.
All twenty-nine pairs pass separate pointwise torus and
meridian equations and finite group enumeration. Six altered
stabilizer, transitivity, meridian and upper-trace inputs reject.

\paragraph{Assurance.}
The uniform graph and torus arguments are written working
proofs, unformalized and externally unreviewed. The exact
negative claim is limited to the specified two degrees and
sixteen diagrams. Other finite images, original double-cover
ranks, worldwide priority and full classical MRC remain open.

\begin{thebibliography}{9}
\small
\bibitem{cornwell} C. Cornwell, Knot contact homology and
representations of knot groups, J. Topology 7 (2014),
1221--1242, Section 5.1.
\url{https://arxiv.org/abs/1303.4943}.
\bibitem{rz} Rost--Zieschang, Meridional generators and plat
presentations of torus links, J. London Math. Soc. (1987).
\url{https://doi.org/10.1112/jlms/s2-35.3.551}.
\bibitem{atlas} ATLAS of Finite Group Representations,
Alternating group $A_8$, Linear group $L_4(2)$.
\url{https://brauer.maths.qmul.ac.uk/Atlas/v3/alt/A8/}.
Primary group order and representation list checked; class
counts above are proved and independently enumerated here.
\bibitem{bk} S. Baader and A. Kjuchukova, Symmetric quotients of
knot groups and a filtration of the Gordian graph (2020).
\url{https://arxiv.org/abs/1711.08144}.
\bibitem{jp} J. Joseph and P. Pongtanapaisan, Meridional rank
and bridge number of knotted 2-spheres, Canad. J. Math. (2025).
Classical finite-cycle construction in preprint Section 3:
\url{https://arxiv.org/html/2111.09233v2}.
\end{thebibliography}
\end{document}
